\section{引言与相关工作}
动态随机存取存储器（DRAM）依靠存储电容的电荷表示信息。访问晶体管导通后，存储节点与位线发生电荷共享，读出系统依据小幅电压差判断存储状态\cite{keeth2001}。粒子入射在硅中产生电子空穴对，随后发生漂移、扩散和复合；生成电荷、端口收集电荷与存储节点净变化具有不同含义\cite{label1996}。因此，粒子使存储电压变化多少，与给定读出过程能否保持正确差分方向，是两个相关但不能等同的问题。

Geant4与TCAD分别适于描述粒子输运和器件电学响应。陈善强、师立勤已将两类工具用于SRAM单粒子效应评估\cite{chen2011}，因此本研究不将工具耦合本身作为首创贡献。机器学习辅助的TCAD单粒子瞬态回归也已有研究\cite{xu2021}。此外，比较重离子、聚焦X射线和脉冲激光诱导响应的研究指出，生成电荷总量相近并不保证器件瞬态一致，空间分布仍影响响应\cite{ryder2021}。这要求联合模型同时说明源项简化、评价对象与验证边界，而非仅给出软件连接和拟合曲线。

本研究以课程提出的AI辅助多特征尺寸存储器评估任务为背景\cite{course}，聚焦三个问题：第一，有限硅区域内的粒子沉积如何转换为可审计的二维TCAD激励；第二，怎样区分无粒子读出偏移与粒子引起的增量；第三，在源项和无粒子基线已知时，低维回归能否描述不同尺寸、粒子和方向下的连续响应。采用简化1T1C单元保留电荷共享与器件输运，暂不模拟完整阵列、真实感放锁存和恢复。

本文将受约束输入生成、等效源检查、匹配评价、预测冻结和全量验证组织为可追溯的评估链。定量结果以统一R50主版本的64独立事件、128响应为主，历史D1与100 Ω诊断只解释设计演进，不并入主矩阵。独立尺寸留出、外推及氧化层和边界验证均按既定接口接入本次封存证据。模型仍未校准真实工艺，本文讨论条件仿真响应和冻结预测误差，而不推断商业DRAM失效率。
